\documentclass[12pt,a4paper]{report}
\usepackage[margin=2.5cm]{geometry}
\usepackage{amsmath}
\usepackage{booktabs}
\usepackage{tabularx}
\usepackage{listings}
\usepackage{xcolor}
\usepackage{titlesec}
\usepackage{parskip}
\usepackage{microtype}
\usepackage[hidelinks]{hyperref}

\definecolor{codebg}{RGB}{245,245,245}
\lstset{basicstyle=\ttfamily\small,backgroundcolor=\color{codebg},frame=single,rulecolor=\color{gray!40},breaklines=true,columns=fullflexible}

% Compact chapter headings suit a shorter report.
\titleformat{\chapter}{\normalfont\LARGE\bfseries}{\thechapter.}{0.6em}{}
\titlespacing*{\chapter}{0pt}{0pt}{1.5em}

\begin{document}

\begin{titlepage}
\centering
\vspace*{1cm}
{\large BSc (Hons) Computer Science\par}
{\large Final Year Project Report\par}
\vspace{3cm}
{\huge\bfseries StudySync: A Shared Revision Planner for University Students\par}
\vspace{3cm}
\begin{tabularx}{0.7\linewidth}{@{}lX@{}}
Student & Your Name \\
Student ID & 12345678 \\
Supervisor & Dr Supervisor Name \\
Second marker & Dr Marker Name \\
Word count & 9\,850 \\
\end{tabularx}
\vfill
{Department of Computing, Example University\par}
{April 2026\par}
\end{titlepage}

\pagenumbering{roman}
\chapter*{Abstract}
In under 300 words, say what you built or investigated, how you evaluated it and what you found.

\chapter*{Statement of originality}
I confirm that this report is my own work and that all sources are acknowledged. \\[2em]
Signature: \rule{5cm}{0.4pt}

\tableofcontents

\clearpage
\pagenumbering{arabic}

\chapter{Introduction}
\section{Problem statement}
Describe the problem your project tackles and who has it.
\section{Aims and objectives}
\begin{enumerate}
  \item Build a working prototype that lets students share revision timetables.
  \item Evaluate usability with at least ten participants.
\end{enumerate}

\chapter{Background research}
Review existing products and relevant academic work.

\chapter{Requirements}
\begin{tabularx}{\linewidth}{@{}llX@{}}
\toprule
ID & Priority & Requirement \\
\midrule
FR1 & Must & Users can create a revision session with a date and subject. \\
FR2 & Must & Users can invite classmates to a shared plan. \\
FR3 & Should & The app sends a reminder before each session. \\
NFR1 & Must & Pages load in under two seconds on a mobile connection. \\
\bottomrule
\end{tabularx}

\chapter{Design and implementation}
\section{Architecture}
Explain the main components and how they talk to each other.
\section{Key code}
Include only short, interesting excerpts.
\begin{lstlisting}[language=Python]
def overlapping(a, b):
    return a.start < b.end and b.start < a.end
\end{lstlisting}

\chapter{Testing and evaluation}
\begin{table}[ht]
\centering
\caption{System Usability Scale results (example).}
\begin{tabular}{lc}
\toprule
Measure & Score \\
\midrule
Mean SUS score & 78.5 \\
Participants & 12 \\
Tasks completed without help & 92\% \\
\bottomrule
\end{tabular}
\end{table}

\chapter{Conclusion}
\section{Achievements}
Relate the outcome to each objective.
\section{Reflection}
Say what went well, what you would change and what you learned.
\section{Future work}
Suggest realistic next steps.

\begin{thebibliography}{9}
\bibitem{nielsen} J.~Author, \emph{Usability Engineering Example}, Publisher, 1993.
\bibitem{sus} K.~Author, A quick usability scale, in \emph{Usability Evaluation in Practice}, 1996, pp.~189--194.
\end{thebibliography}

\appendix
\chapter{Project plan}
Include your original Gantt chart or plan here.
\chapter{User guide}
Explain how to install and run the software.

\end{document}
